\section{Depth-First Search (DFS) Traversal}
\label{sec:graph-dfs}

\problemstatement{Given a graph (as an adjacency list) and a starting
node, visit every node reachable from the start exactly once, exploring
as deep as possible along each path before backtracking.}

\begin{patternbox}
\emph{``Visit every reachable node, going deep before wide''} is solved
by recursion (or an explicit stack): visit the current node, mark it
visited, then recurse into every unvisited neighbor in turn.
\end{patternbox}

\subsection*{Algorithm}

\begin{enumerate}
  \item Maintain a $\textit{visited}$ array, initialized to all false.
  \item Define $\textit{dfs}(\textit{node})$: mark $\textit{node}$ as
        visited, add it to the result; for every neighbor of
        $\textit{node}$ not yet visited, recursively call
        $\textit{dfs}(\textit{neighbor})$.
  \item To traverse the whole graph (which may be disconnected), call
        $\textit{dfs}(i)$ for every node $i$ not yet visited.
\end{enumerate}

\subsection*{Why a \textit{visited} array is required}

Without tracking visited nodes, a cycle in the graph (or even a simple
back-and-forth edge between two nodes) would cause infinite recursion.
Marking a node visited \textbf{before} recursing into its neighbors
(not after returning) is essential: it guarantees that if two different
neighbors both point back to an already-in-progress node, only the
first one to reach it recurses into it.

\subsection*{C++ implementation}

\begin{lstlisting}
#include <bits/stdc++.h>
using namespace std;

void dfs(int node, vector<vector<int>>& adj, vector<bool>& visited,
         vector<int>& result) {
    visited[node] = true;
    result.push_back(node);

    for (int neighbor : adj[node]) {
        if (!visited[neighbor]) {
            dfs(neighbor, adj, visited, result);
        }
    }
}

vector<int> dfsTraversal(int n, vector<vector<int>>& adj) {
    vector<bool> visited(n, false);
    vector<int> result;

    for (int i = 0; i < n; i++) {
        if (!visited[i]) {
            dfs(i, adj, visited, result);
        }
    }
    return result;
}

int main() {
    int n = 5;
    vector<vector<int>> adj(n);
    auto addEdge = [&](int u, int v) { adj[u].push_back(v); adj[v].push_back(u); };
    addEdge(0, 1); addEdge(0, 2); addEdge(1, 3); addEdge(2, 4);

    for (int x : dfsTraversal(n, adj)) cout << x << " ";
    cout << endl; // 0 1 3 2 4
    return 0;
}
\end{lstlisting}

\subsection*{Dry run}

\dryrun{Take the graph with edges $0{-}1$, $0{-}2$, $1{-}3$, $2{-}4$,
starting traversal from node $0$.}

Visit $0$ (mark visited, record). Its neighbors: $1$ (unvisited) --
recurse. Visit $1$, record. Its neighbors: $0$ (visited, skip), $3$
(unvisited) -- recurse. Visit $3$, record; no unvisited neighbors,
return. Back at $1$: no more neighbors, return. Back at $0$: next
neighbor $2$ (unvisited) -- recurse. Visit $2$, record. Its neighbors:
$0$ (visited, skip), $4$ (unvisited) -- recurse. Visit $4$, record; no
unvisited neighbors, return. Final order: $0,1,3,2,4$.

\begin{complexitybox}
\textbf{Time Complexity.} $O(V+E)$ -- every vertex is visited once
(marking it visited), and every edge is examined once (from each
endpoint, checking if the other side is visited).
\textbf{Space Complexity.} $O(V)$ for the \textit{visited} array and
the recursion stack (which can be as deep as $V$ in a path-shaped
graph).
\end{complexitybox}

\begin{takeawaybox}
DFS's ``visit, mark, recurse into unvisited neighbors'' template is the
single most reused piece of code in this chapter -- nearly every later
problem either calls this function directly or modifies one small part
of it (the visiting action, the stopping condition, or what gets
collected).
\end{takeawaybox}
